\documentclass[11pt,letterpaper]{article}

\usepackage[top=1in, bottom=1in, left=1in, right=1in, headheight=14pt]{geometry}
\usepackage{fontspec}
\setmainfont{DejaVu Serif}
\setsansfont{DejaVu Sans}
\setmonofont{DejaVu Sans Mono}

\usepackage{xcolor}
\usepackage{titlesec}
\usepackage{fancyhdr}
\usepackage{enumitem}
\usepackage{hyperref}
\usepackage{booktabs}
\usepackage{tabularx}
\usepackage{longtable}
\usepackage{colortbl}
\usepackage{multirow}
\usepackage{array}
\usepackage{parskip}
\usepackage{tcolorbox}
\usepackage{graphicx}
\usepackage{lastpage}

%% COLOR SCHEME — Broadcast / Signal domain
%% Primary: deep indigo (carrier wave, transmission depth)
%% Secondary: phosphor teal (CRT glow, signal green)
%% Tertiary: hot amber (analog warmth, phosphor orange)

\definecolor{primary}{HTML}{1A237E}
\definecolor{secondary}{HTML}{006064}
\definecolor{tertiary}{HTML}{BF360C}
\definecolor{accent}{HTML}{283593}
\definecolor{lightprimary}{HTML}{E8EAF6}
\definecolor{lightsecondary}{HTML}{E0F2F1}
\definecolor{lighttertiary}{HTML}{FBE9E7}
\definecolor{rowgray}{HTML}{F5F5F5}
\definecolor{bordergray}{HTML}{BDBDBD}
\definecolor{subtlegray}{HTML}{757575}
\definecolor{darktext}{HTML}{212121}

%% SECTION FORMATTING
\titleformat{\section}
  {\Large\bfseries\sffamily\color{primary}}
  {\thesection}{1em}{}
  [\vspace{-0.5em}\textcolor{bordergray}{\rule{\textwidth}{0.4pt}}]

\titleformat{\subsection}
  {\large\bfseries\sffamily\color{secondary}}
  {\thesubsection}{1em}{}

\titleformat{\subsubsection}
  {\normalsize\bfseries\sffamily\color{tertiary}}
  {\thesubsubsection}{1em}{}

\titlespacing*{\section}{0pt}{1.5em}{0.8em}
\titlespacing*{\subsection}{0pt}{1.2em}{0.5em}
\titlespacing*{\subsubsection}{0pt}{0.8em}{0.3em}

%% CUSTOM ENVIRONMENTS
\newcommand{\stageheader}[2]{%
  \noindent\colorbox{#2}{\makebox[\textwidth][l]{%
    \textcolor{white}{\bfseries\sffamily\hspace{6pt}#1\hspace{6pt}}}}%
  \vspace{0.5em}\par%
}

\newtcolorbox{asciibox}{
  colback=rowgray, colframe=bordergray,
  arc=1pt, boxrule=0.5pt,
  left=6pt, right=6pt, top=4pt, bottom=4pt,
  fontupper=\small\ttfamily,
  breakable
}

\newtcolorbox{quotebox}{
  colback=lightprimary, colframe=primary,
  arc=2pt, boxrule=0.5pt,
  left=12pt, right=12pt, top=8pt, bottom=8pt,
  breakable
}

\newcommand{\metafield}[2]{\textbf{\sffamily #1} #2\par}

%% TABLE HELPERS
\newcolumntype{L}[1]{>{\raggedright\arraybackslash}p{#1}}
\newcolumntype{C}[1]{>{\centering\arraybackslash}p{#1}}
\newcommand{\tableheader}[1]{\textbf{\sffamily\textcolor{white}{#1}}}

%% HEADERS / FOOTERS
\pagestyle{fancy}
\fancyhf{}
\fancyhead[L]{\small\sffamily\textcolor{subtlegray}{The Polyrhythm Standard}}
\fancyhead[R]{\small\sffamily\textcolor{subtlegray}{GSD Mission Package}}
\fancyfoot[C]{\small\sffamily\textcolor{subtlegray}{Page \thepage\ of \pageref{LastPage}}}
\renewcommand{\headrulewidth}{0.4pt}
\renewcommand{\footrulewidth}{0pt}

%% HYPERLINKS
\hypersetup{
  colorlinks=true,
  linkcolor=secondary,
  urlcolor=secondary,
  pdfauthor={GSD Ecosystem / Tibsfox},
  pdftitle={The Polyrhythm Standard -- Mission Package},
  pdfsubject={Vision to Mission Pipeline},
}

%% DOCUMENT
\begin{document}

%% =====================================================================
%% TITLE PAGE
%% =====================================================================
\thispagestyle{empty}
\begin{center}
\vspace*{2.5cm}

{\small\sffamily\textcolor{primary}{\textbf{GSD RESEARCH MISSION PACKAGE}}}

\vspace{0.3cm}
\textcolor{bordergray}{\rule{8cm}{0.4pt}}

\vspace{1.5cm}
{\fontsize{26}{32}\selectfont\bfseries\sffamily\textcolor{primary}{The Polyrhythm Standard\\[0.3em]}}
{\fontsize{16}{20}\selectfont\bfseries\sffamily\textcolor{secondary}{50/60\,Hz, NTSC/PAL, SMPTE, High-Frequency Displays,\\[0.2em]8K Video, and Large-Format Digital Imaging}}

\vspace{0.8cm}
{\Large\sffamily\textcolor{tertiary}{Signal Engineering \& Visual Time Architecture}}

\vspace{0.5cm}
{\large\sffamily\itshape\textcolor{subtlegray}{A Deep Research Study}}

\vspace{3cm}
{\sffamily\textcolor{accent}{Vision $\rightarrow$ Research $\rightarrow$ Mission}}\\[0.3em]
{\small\sffamily\textcolor{accent}{Full Pipeline Execution}}

\vspace{2.5cm}
{\small\sffamily\textcolor{subtlegray}{Date: March 26, 2026 \quad|\quad Status: Ready for GSD Orchestrator}}\\[0.3em]
{\small\sffamily\textcolor{subtlegray}{Activation Profile: Fleet \quad|\quad Est.\ 8--10 context windows across 4--5 sessions}}
\end{center}
\newpage

%% =====================================================================
%% TABLE OF CONTENTS
%% =====================================================================
\tableofcontents
\newpage

%% =====================================================================
%% STAGE 1 — VISION DOCUMENT
%% =====================================================================
\stageheader{STAGE 1 --- VISION DOCUMENT}{primary}

\section{The Polyrhythm Standard --- Vision Guide}

\metafield{Date:}{March 26, 2026}
\metafield{Status:}{Initial Vision / Cold-Start Research Complete}
\metafield{Depends On:}{gsd-skill-creator v1.49+, DACP milestone}
\metafield{Context:}{Six-module Fleet mission; trilinear CCD and SMPTE timecode as architectural exemplars}

\subsection{Vision}

There is a rhythm at the heart of every screen you have ever watched, embedded so deeply it predates color
television. When electrical engineers in the 1940s and 1950s wired the grids of North America at
60\,Hz and Europe at 50\,Hz, they made a geopolitical decision that would reverberate through every
broadcast, every edit bay, every game console, every cinema frame, and every museum-grade digitization
workflow for the next eight decades. The television picture rate was locked to the local power frequency
to avoid beating artifacts with fluorescent lighting and power hum. The Atlantic Ocean thus became a
temporal fault line: 30 frames per second on one side, 25 on the other.

These two frequencies --- 50\,Hz and 60\,Hz --- do not share a common sub-multiple below 10. Their
Least Common Multiple is 300\,Hz. This is not a footnote. It is the hidden rhythm beneath every
standards-conversion device, every PAL-to-NTSC frame-rate shimmy, every 3:2 pulldown that stretches
24-fps film to 59.94 fields, and every modern high-frequency display that resolves the tension by
simply running fast enough for both to divide evenly into it. The 50/60 polyrhythm is the oldest
living constraint in consumer electronics, and the entire history of video engineering is, in some
sense, the story of building machines that can hear both rhythms at once.

SMPTE timecode is the first formal language invented to count those frames, and its drop-frame
mechanism is a confession: the color NTSC system's frame rate is not 30 frames per second at all.
When engineers added color to the monochrome signal in 1953, they needed the color subcarrier
(3.579545\,MHz, precisely 315/88\,MHz) to interleave with the luminance harmonics without visible
beating. The only clean solution was to slow every timing parameter by a factor of 1000/1001 --- a
0.1\% reduction. The frame rate became 29.970030\,Hz. That invisible 0.1\% gap accumulates into 3.6
seconds of drift per broadcast hour, 86.4 seconds of drift per broadcast day. Drop-frame timecode
skips two frame numbers at the start of every minute (except every tenth minute) to stay synchronized
with the wall clock. It is a rhythmic correction, not unlike a musician adding a grace note at the
end of a phrase to arrive at the downbeat on time.

High-frequency displays represent the upstream end of this lineage. Modern 120\,Hz displays emerged
partly because 120 is the LCM of 24, 25, and 30 --- it can divide cleanly into film-at-24, PAL-at-25,
and NTSC-at-30 without judder. The push to 240\,Hz, 360\,Hz, 500\,Hz, and beyond is driven by human
physiology: the motion processing system of the visual cortex continues to extract useful information
at frequencies well beyond 60\,Hz, and sample-and-hold displays produce motion blur that only resolves
as frame time decreases. The resolution frontier is not resolution at all --- it is temporal density.

At the high end of still imaging, the scanning back tells the same story in spatial rather than
temporal terms. Where a single-shot matrix sensor captures a Bayer-filtered interpolation of color,
the trilinear CCD scanning back makes three separate passes through red, green, and blue for every
single pixel column, then scans the entire sensor width mechanically. This is the photographic
equivalent of drop-frame correction: accepting slowness in exchange for absolute accuracy. Phase One's
150-megapixel IQ4 back and BetterLight's Super8K system do not compete with matrix sensors on speed;
they compete on truth. The spaces between their measurements are precisely calibrated.

This mission documents the engineering archaeology of these standards, traces the polyrhythm from
power grid to pixel, and builds a GSD research package suitable for multi-module Fleet execution.

\subsection{Problem Statement}

\begin{enumerate}[leftmargin=*, label=\textbf{\arabic*.}]
  \item \textbf{The 50/60 polyrhythm remains unresolved in modern workflows.}
  Content produced in PAL (25\,fps) environments must cross into NTSC (29.97\,fps) broadcast chains
  and vice versa. Standards conversion introduces either judder, motion blur, or temporal artifacts
  because the two rates share no common sub-multiple below 300. Modern IP-based systems have
  dissolved many format barriers but the underlying temporal mismatch persists in every live
  broadcast, sports event, and theatrical exhibition.

  \item \textbf{SMPTE timecode's drop-frame mechanism is routinely misunderstood, causing sync failures.}
  The 0.1\% frequency offset between ``30\,fps'' and true 29.97\,fps is a mathematical constant that
  propagates into DAWs, NLEs, and broadcast playout systems. Editors who mix drop-frame and
  non-drop-frame timecode within a single project introduce 3.6 seconds of drift per hour --- enough
  to miss broadcast windows, misalign audio, and corrupt archival timestamps.

  \item \textbf{High-frequency displays have outpaced the understanding of their underlying science.}
  Marketing has conflated ``motion rate,'' ``refresh rate,'' ``response time,'' and ``input lag'' into
  a single undifferentiated metric. The actual perceptual science --- MPRT (Moving Picture Response
  Time) versus GtG (gray-to-gray), sample-and-hold motion blur, VRR, and black frame insertion --- is
  poorly communicated to both buyers and developers.

  \item \textbf{8K and ultra-HD standards are fragmented across competing ecosystems.}
  ITU-R BT.2020, the 8K Association, LG/CTA's competing standard, NHK's Super Hi-Vision (SHV), and
  ATSC 3.0 broadcasting each define slightly different capability sets. HDR sub-standards
  (HDR10, HDR10+, HLG, Dolby Vision) add a second axis of fragmentation. As Samsung holds near-monopoly
  over 8K consumer television and Dolby Vision exits the segment, the content-display contract for
  8K is unresolved.

  \item \textbf{Scanning backs and ultra-high-resolution still imaging lack a unified technical reference.}
  The trilinear CCD scanning back --- the closest digital equivalent to large-format sheet film --- is
  undersupported in educational and practitioner literature. BetterLight's Super8K, Phase One's IQ4 150MP,
  and Hasselblad's multi-shot 400MP compositing mode each represent distinct resolution strategies that
  are poorly contextualized against each other and against 8K cinema acquisition.
\end{enumerate}

\subsection{Core Concept}

\textbf{The Polyrhythm Standard:} two incommensurable temporal grids (50\,Hz / 60\,Hz) originating in
infrastructure, propagating into every display, broadcast, and imaging standard, and resolved only by
running fast enough (300\,Hz LCM), accurate enough (drop-frame), or slow enough (scanning back) to
honor both.

\subsubsection{Domain Map --- The Temporal Stack}

\begin{asciibox}
  SPATIAL RESOLUTION                         TEMPORAL RESOLUTION
  (still imaging)                            (motion imaging)
  
  Scanning Back (trilinear CCD, 3-pass)      1000Hz+ displays (esports)
    |  BetterLight Super8K: 12000x15990       |  360-500Hz competitive
    |  Phase One IQ4: 150MP native            |  240Hz premium consumer
    |  Hasselblad H6D-400c MS: 400MP          |  120Hz: LCM of 24/25/30
    v  multi-shot composite                   v  60Hz: NTSC field rate
  Medium Format (single-shot area array)        50Hz: PAL field rate
    |  100-150MP CMOS matrix                SMPTE Drop Frame
    |  Bayer filter / color interpolation     |  29.97 = 30/1.001 fps
    v  Phase One, Hasselblad, Fuji GFX        |  2 frame-numbers dropped
  35mm Full Frame                               |  per minute, not every 10th
    |  24-61MP CMOS                           v  86.4ms residual drift/day
    v  Canon, Sony, Nikon                   NTSC / PAL Standards
  Cinema 4K-8K                               |  525 lines 60 fields (NTSC)
    |  REDcode RAW, XAVC, ProRes              |  625 lines 50 fields (PAL)
    v  7680x4320 @ 120fps emerging            v  Power grid origin: 60/50Hz AC
  Consumer 4K → 8K                         Electrical Infrastructure
    v  ITU-R BT.2020, HEVC/AV1              |  North America / Japan: 60Hz
  Broadcast Delivery                         |  Europe / Africa / Asia: 50Hz
    v  ATSC 3.0, DVB-T2, NHK SHV 8K        v  Historical fork: 1940s-1950s
\end{asciibox}

\subsubsection{Module Architecture --- Six Research Modules}

\begin{asciibox}
  [A] Temporal Origins         [D] The Polyrhythm Architecture
      50/60Hz power grid            LCM=300Hz, standards conversion
      NTSC/PAL divergence           3:2 pulldown, Tetra-sync
      SECAM third standard          PAL-to-NTSC artifacts
              |                              |
              v                              v
  [B] NTSC & PAL Engineering   [E] High-Frequency Refresh Science
      Color subcarrier              120Hz LCM, 240-1000Hz arc
      525/625 line systems          MPRT vs GtG, sample-and-hold
      Interlace vs progressive      VRR / FreeSync / G-Sync / BFI
              |                              |
              v                              v
  [C] SMPTE Timecode           [F] Large Format & Scanning Backs
      29.97 drift anatomy           8K/SHV/ITU-R BT.2020
      Drop-frame correction         Medium format area sensors
      LTC / VITC / AES3             Trilinear CCD scanning backs
      Broadcast sync chains         Multi-shot compositing

  Modules A+B run in parallel (Wave 1 Track A)
  Modules C+D run in parallel (Wave 1 Track B)
  Modules E+F run in parallel (Wave 1 Track C)
  Wave 2: cross-module synthesis (polyrhythm through-line)
\end{asciibox}

\subsection{Research Modules}

\subsubsection{Module A --- Temporal Origins}
Traces the 50/60\,Hz power-grid fork from 1940s infrastructure decisions through the establishment
of NTSC (1941/1953), PAL (1963), and SECAM (1956). Documents why picture rate was locked to power
frequency, the interlace bandwidth compromise, and the three-standard world it created. Geographic
coverage includes North America, Japan, Western Europe, Eastern Europe/France, and the developing
world's adoption patterns.

\subsubsection{Module B --- NTSC and PAL Engineering}
Technical dissection of both color encoding systems: NTSC's quadrature amplitude modulation of
I and Q signals onto a 3.579545\,MHz subcarrier; PAL's phase-alternating line trick that
self-corrects for phase errors; the Hanover bars artifact and PAL's 1-line delay correction circuit.
Comparison of 525-line/59.94\,Hz versus 625-line/50\,Hz parameters including horizontal sync,
blanking intervals, vertical resolution, and color bandwidth.

\subsubsection{Module C --- SMPTE Timecode and the Drop-Frame Problem}
Origin of the 0.1\% frequency offset in NTSC color standardization. Drop-frame timecode as a
mathematical correction: 18 frame-numbers dropped per 10-minute block, residual error of 1.0\,ppm
(86.4\,ms/day). LTC (Longitudinal Timecode), VITC (Vertical Interval Timecode), and AES3 timecode
formats. Modern ST\,12-3 extensions for frame rates above 40\,fps, up to 120\,fps. Broadcast sync
chains: black burst, tri-level sync, PTP (IEEE 1588).

\subsubsection{Module D --- The Polyrhythm Architecture}
LCM(50, 60) = 300\,Hz as the mathematical resolution of the 50/60 tension. Standards conversion
methods: 3:2 pulldown (24fps film to NTSC 59.94 fields), frame-rate conversion artifacts (motion blur,
judder), motion-compensated interpolation. The role of 120\,Hz as LCM(24, 25, 30) in modern TV.
IP-era convergence: does baseband signal have a future? HFR cinema (48fps, 60fps, 120fps).

\subsubsection{Module E --- High-Frequency Refresh Displays}
Physics of sample-and-hold motion blur on LCD/OLED. The MPRT/GtG distinction. Scientific literature
on human visual motion perception: PMC 8777290 (Frontiers in Neuroscience, 2022) recommending
$\geq$120\,Hz for motion studies, $\geq$240\,Hz for high-velocity stimuli. The geometric refresh
progression: 60\,$\rightarrow$\,120\,$\rightarrow$\,240\,$\rightarrow$\,480\,$\rightarrow$\,960\,Hz.
Current state: 360-500\,Hz production panels (Dell Alienware AW2523HF). Black frame insertion (BFI),
Variable Refresh Rate (VRR) standards. Input latency: 28\,ms system latency reduction at 360\,Hz vs
60\,Hz in measured Fortnite benchmarks.

\subsubsection{Module F --- Large Format, 8K, and Scanning Backs}
Ultra-high-resolution still imaging: Phase One IQ4 150MP (53.7\,$\times$\,40.4\,mm sensor),
Hasselblad X2D II 100C (43.8\,$\times$\,32.9\,mm), Fuji GFX 100S II. Multi-shot pixel-shift:
Hasselblad H6D-400c MS (400MP composite). Scanning backs: BetterLight Super8K (12,000\,$\times$\,15,990
pixels, 1.1\,GB 48-bit TIFF, trilinear CCD). Comparison to 8K cinema (7680\,$\times$\,4320 = 33.2\,MP).
8K standards: ITU-R BT.2020, SMPTE ST 2036-1, NHK Super Hi-Vision (launched December 2018). HDR
standards: HDR10, HDR10+, HLG, Dolby Vision; post-2026 fragmentation with Samsung near-monopoly
at consumer 8K.

\subsection{Chipset Configuration}

\begin{asciibox}
chipset:
  agnus:            # Orchestration / DMA / context management
    model: opus
    role: FLIGHT + PLAN + INTEG + RETRO
    focus:
      - Cross-module polyrhythm synthesis
      - Standards genealogy analysis
      - Cultural and technical sensitivity
      - Temporal accuracy (timecode, frame-rate math)

  denise:           # Display / output rendering
    model: sonnet
    role: EXEC x3 + VERIFY + LOG
    focus:
      - Module document production (A through F)
      - Source validation and citation accuracy
      - Wave 3 publication assembly

  paula:            # Audio / I/O / anticipatory
    model: haiku
    role: SCOUT + BUDGET + shared schemas
    focus:
      - Shared schema generation (frequency tables, frame-rate matrices)
      - Token budget tracking
      - Source index scaffolding

  gary_68000:       # CPU / glue logic
    model: sonnet
    role: CRAFT-SIGNAL + CRAFT-OPTICS
    focus:
      - CRAFT-SIGNAL: broadcast engineering, SMPTE, sync standards
      - CRAFT-OPTICS: large-format photography, sensor architecture
\end{asciibox}

\subsection{Scope Boundaries}

\subsubsection{In Scope (v1.0)}
\begin{itemize}[leftmargin=*]
  \item 50/60\,Hz power-grid origin and NTSC/PAL/SECAM divergence history
  \item SMPTE 29.97 drop-frame mathematics and all major timecode formats (LTC, VITC, AES3)
  \item High-frequency refresh display science from 60\,Hz through current 500\,Hz production hardware
  \item 8K consumer/broadcast standards (ITU-R BT.2020, 8K Association, NHK SHV, ATSC 3.0)
  \item HDR standards: HDR10, HDR10+, HLG, Dolby Vision --- current fragmentation state
  \item Medium format digital photography (Phase One, Hasselblad, Fujifilm GFX)
  \item Trilinear CCD scanning backs (BetterLight, Kigamo, Anagramm)
  \item Multi-shot / pixel-shift compositing (Hasselblad H6D-400c MS)
  \item Standards conversion artifacts (3:2 pulldown, motion-compensated interpolation)
\end{itemize}

\subsubsection{Out of Scope (v1.0)}
\begin{itemize}[leftmargin=*]
  \item SECAM engineering deep-dive (mentioned for completeness, not detailed)
  \item Audio timecode formats (MTC, MIDI clock) beyond broadcast context
  \item Virtual Reality and XR display systems
  \item Cinema projector standards (DCI, Dolby Cinema) beyond reference
  \item Satellite broadcast delivery infrastructure
  \item Film scanning and photochemical workflows
\end{itemize}

\subsection{Success Criteria}

\begin{enumerate}[leftmargin=*, label=\textbf{SC-\arabic*.}]
  \item Module A documents the power-grid to picture-rate causal chain with primary historical sources
        covering NTSC (1941/1953), PAL (1963), and geographic adoption patterns.
  \item Module B provides complete engineering parameters for both NTSC 525/59.94 and PAL 625/50
        systems including color subcarrier frequencies, line timing, and color correction mechanisms.
  \item Module C traces the 0.1\% frequency offset from its color-compatibility origin through the
        drop-frame correction mechanism with exact mathematical values and residual error.
  \item Module C documents all major SMPTE timecode variants (LTC, VITC, ST 12-3 HFR extensions)
        with bit-field layouts and use cases.
  \item Module D establishes LCM(50, 60) = 300\,Hz as a unifying structural concept and traces its
        implications through standards conversion, modern TV, and HFR cinema.
  \item Module D documents at least three standards conversion methods (3:2 pulldown, frame-rate
        conversion interpolation, and motion-compensated conversion) with artifact descriptions.
  \item Module E documents the scientific evidence base for display refresh rates at 60, 120, 240,
        and 360+\,Hz including at least one peer-reviewed neuroscience citation.
  \item Module E distinguishes MPRT from GtG, explains sample-and-hold motion blur physics, and
        describes VRR, BFI, and strobed backlight operation.
  \item Module F documents the complete medium format sensor landscape (Phase One IQ4, Hasselblad
        X2D II, Fuji GFX 100S II) with sensor dimensions, bit depth, and dynamic range.
  \item Module F documents trilinear CCD scanning back operation, key commercial systems (BetterLight
        Super8K), and cultural heritage archival applications with file size data.
  \item Module F documents 8K standards (ITU-R BT.2020, NHK SHV, SMPTE ST 2036-1) and current
        consumer 8K market state including HDR standard fragmentation.
  \item The synthesis module traces the polyrhythm concept coherently from power grid through
        NTSC/PAL divergence to SMPTE drift to LCM display resolution to scanning-back temporal patience.
\end{enumerate}

\subsection{Relationship to Other Documents}

\begin{center}
\rowcolors{2}{rowgray}{white}
\begin{tabularx}{\textwidth}{L{4cm}XC{2.5cm}}
\toprule
\rowcolor{primary}
\tableheader{Document} & \tableheader{Relationship} & \tableheader{Direction} \\
\midrule
GSD Deep Audio Analyzer Mission & Overlaps on SMPTE sync, sample rate relationships & Sibling \\
Fox Infrastructure Group Plan & Broadcast infrastructure cross-domain (FoxFiber, FoxData) & Downstream \\
The Space Between (textbook) & Mathematical foundations: LCM as set-theory intersection concept & Upstream \\
GSD Chipset Architecture Vision & Agnus/Denise/Paula model maps onto audio/video pipeline metaphor & Parallel \\
College of Knowledge & Signal standards module candidate; cross-language concept translation & Downstream \\
\bottomrule
\end{tabularx}
\end{center}

\subsection{The Through-Line}

The history of visual standards is the history of spaces between. Between 50 and 60 lives the
LCM at 300. Between 30 and 29.97 lives the 3.6-second hourly drift that drop-frame corrects with
rhythmic grace. Between the single-shot Bayer matrix and the film-print lives the scanning back's
trilinear patience, measuring red before green before blue, finding color in the pauses between passes.

The Amiga Principle applies here at a civilizational scale: the 50 and 60\,Hz standards were not
mistakes. They were the most efficient encodings of local physical constraint --- the power frequency
beneath every CRT, the phosphor decay rate, the bandwidth that the electromagnetic spectrum could carry.
The divergence produced incompatibility, but also diversity: PAL's self-correcting phase alternation
was a more elegant color solution than NTSC's manual hue control; NTSC's higher field rate produced
smoother motion. Neither was simply better. Both were architecturally principled responses to local
conditions. The spaces between them --- the standards-conversion apparatus, the drop-frame workaround,
the 120\,Hz LCM display --- are not failures of coordination. They are the connective tissue of a
genuinely distributed visual world. This mission documents that tissue, understands its structure,
and prepares it for execution.

\newpage

%% =====================================================================
%% STAGE 2 — RESEARCH REFERENCE
%% =====================================================================
\stageheader{STAGE 2 --- RESEARCH REFERENCE}{secondary}

\section{The Polyrhythm Standard --- Research Reference}

\subsection{How to Use This Document}

This research reference provides verified findings for all six modules. Each section contains
quantitative data attributed to specific sources. The bibliography is organized by source category.
All sources are professional organizations, standards bodies, peer-reviewed journals, or manufacturer
technical documentation. No entertainment media or unsourced claims are included.

\textbf{Key source organizations:}
\begin{itemize}[leftmargin=*]
  \item \textbf{SMPTE} (Society of Motion Picture and Television Engineers) --- ST\,12M, ST\,2036-1
  \item \textbf{ITU-R} (International Telecommunication Union, Radiocommunication Sector) --- BT.2020, BT.709
  \item \textbf{IEEE} --- 802.1AS (PTP), synchronization standards
  \item \textbf{NHK} (Nippon Hoso Kyokai / Japan Broadcasting Corporation) --- Super Hi-Vision research
  \item \textbf{EBU} (European Broadcasting Union) --- 25\,fps timecode, PAL broadcast engineering
  \item \textbf{Frontiers in Neuroscience / PMC} --- Peer-reviewed motion perception research
  \item \textbf{Phase One, Hasselblad, BetterLight} --- Manufacturer technical specifications
  \item \textbf{CTA} (Consumer Technology Association) --- 8K TV standards
  \item \textbf{8K Association} --- Industry 8K specification coalition
\end{itemize}

\subsection{Module A Reference --- Temporal Origins}

\subsubsection{Power Grid to Picture Rate}

The causal chain from electrical infrastructure to video standard is direct and documented. North
American power grids settled on 60\,Hz AC supply in the early twentieth century, while European grids
standardized on 50\,Hz. Early television engineers chose to synchronize the vertical field rate to
the local power frequency to prevent visible interference patterns (``hum bars'') between the scanning
electron beam and electromagnetic fields from power wiring (The Broadcast Bridge, 2022).

The original U.S. monochrome television standard specified an exact 60.000\,Hz field rate, yielding
exactly 30 frames per second in interlaced format (525 lines per frame, 262.5 lines per field). The
EBU standard for Europe specified 50\,Hz field rate, 25 frames per second, 625 lines per frame
(312.5 lines per field). The difference in line counts (525 vs.\ 625) reflected different bandwidth
allocations: more lines required more horizontal bandwidth, and European standards were able to
allocate more bandwidth per channel.

\subsubsection{The Three-Standard World}

\begin{center}
\rowcolors{2}{rowgray}{white}
\begin{tabularx}{\textwidth}{L{2cm}L{3cm}C{2cm}C{2cm}C{2.5cm}}
\toprule
\rowcolor{secondary}
\tableheader{Standard} & \tableheader{Regions} & \tableheader{Lines} & \tableheader{Field Rate} & \tableheader{Color System} \\
\midrule
NTSC & N.\ America, Japan & 525 & 59.94\,Hz & QAM on 3.58\,MHz \\
PAL  & W.\ Europe, Australia, China, India & 625 & 50\,Hz & Phase-alt on 4.43\,MHz \\
SECAM & France, E.\ Europe, Middle East & 625 & 50\,Hz & FM sequential \\
\bottomrule
\end{tabularx}
\end{center}

NTSC (National Television System Committee) color was standardized by the FCC in December 1953.
PAL (Phase Alternating Line) was developed by Walter Bruch at Telefunken in Hanover, West Germany,
patented in December 1962, and first broadcast in the UK in July 1967 and West Germany in August 1967
(Wikipedia, PAL). SECAM (Séquentiel Couleur à Mémoire) was developed by Henri de France and first
broadcast in France in 1967.

PAL's key engineering advantage: the phase of the color subcarrier is deliberately inverted on
alternate lines. A simple 1-line delay circuit in the receiver allows phase errors to cancel,
making PAL more robust against multipath propagation and antenna alignment errors than NTSC.
This is why NTSC earned the informal nickname ``Never The Same Color.''

\subsection{Module B Reference --- NTSC and PAL Engineering}

\subsubsection{NTSC Color Subcarrier Derivation}

The original NTSC monochrome standard used a 30.000\,Hz frame rate. When color was added in 1953,
the 3.898125\,MHz candidate subcarrier produced visible interference with the 4.500\,MHz audio
intercarrier. Engineers reduced the subcarrier multiple from 495 to 455 (half-line multiples) and
moved the beat frequency from approximately 600\,kHz to approximately 920\,kHz. To make this beat
frequency exactly equal to an odd multiple of half the line scan frequency, they reduced all timing
parameters by 0.1\% (a factor of 1000/1001). The final values:
\begin{itemize}[leftmargin=*]
  \item Color subcarrier: $315/88\,\text{MHz} = 3.57954545\ldots\,\text{MHz}$
  \item Line scan frequency: $9/572\,\text{MHz} = 15734.27\,\text{Hz}$
  \item Frame rate: $30 \times 1000/1001 = 29.97002997\ldots\,\text{fps}$
\end{itemize}
(SMPTE Timecode article, Wikipedia; Phil Rees Timecode Reference)

\subsubsection{PAL Color Subcarrier}

PAL uses a subcarrier of 4.43361875\,MHz (PAL\,B/G/H/I/D). The phase alternation occurs on alternate
lines: the subcarrier is rotated by $+90°$ on odd lines and $-90°$ on even lines relative to the
reference burst. The receiver adds the current line to the 1-line-delayed previous line; phase errors
cancel in the summation while chrominance information (which is $180°$ different between lines) combines
constructively. PAL's 625-line system provides 576 active lines of vertical resolution compared to
NTSC's 480 active lines --- a 20\% resolution advantage.

\subsection{Module C Reference --- SMPTE Timecode}

\subsubsection{The 29.97 Problem}

The 0.1\% reduction in NTSC timing caused a specific problem: timecode that counted frames at
exactly 30 per second (nominal) while the actual frame rate was 29.97\,fps caused cumulative drift.
After one hour, the timecode read 01:00:00:00 while the wall clock showed 01:00:03.6 seconds. Over
a broadcast day, the error approaches 86.4 seconds (SMPTE Timecode, Wikipedia; SMPTE Grokipedia, 2026).

\subsubsection{Drop-Frame Correction}

Drop-frame timecode, formalized in SMPTE 12M-1986, corrects this by skipping frame numbers 0 and 1
at the start of every minute except minutes divisible by 10. Specifically:
\begin{itemize}[leftmargin=*]
  \item 18 frame-numbers dropped per 10-minute block (= 18,000 frames at 30\,fps)
  \item Effective rate: $30 \times 0.999 = 29.97\,\text{fps}$
  \item Residual error: $1.0\,\text{ppm}$ (approximately 86.4\,ms per day)
\end{itemize}
No actual video frames are dropped. Only the numbering labels skip. Drop-frame notation uses
semicolons (;) as separators rather than colons (:).

\subsubsection{Timecode Formats}

\begin{center}
\rowcolors{2}{rowgray}{white}
\begin{tabularx}{\textwidth}{L{2.5cm}XC{2cm}}
\toprule
\rowcolor{secondary}
\tableheader{Format} & \tableheader{Description} & \tableheader{Standard} \\
\midrule
LTC & Longitudinal Timecode: biphase-mark encoded audio signal on dedicated track & ST 12-1 \\
VITC & Vertical Interval Timecode: embedded in video blanking interval, readable at pause & ST 12-1 \\
ATC & Ancillary Timecode: embedded in serial digital video SDI ancillary data space & ST 12-2 \\
HFR TC & High Frame Rate extension for rates above 40\,fps (up to 120\,fps via user bits) & ST 12-3 \\
\bottomrule
\end{tabularx}
\end{center}

SMPTE ST 12 supports frame rates of: 23.976, 24, 25, 29.97, 30, 47.95, 48, 50, 59.94, and 60\,fps,
with ST 12-3 extending support to 120\,fps (SMPTE Grokipedia, 2026).

\subsubsection{Broadcast Sync}

Modern broadcast synchronization uses IEEE\,1588 Precision Time Protocol (PTP), which replaces the
analog black-burst (NTSC) and tri-level sync (HD) reference generators. PTP provides sub-microsecond
accuracy over IP networks, enabling synchronized multi-camera production without physical sync cables.
This represents the convergence point where the 50/60\,Hz physical reference becomes a digital
timestamp.

\subsection{Module D Reference --- The Polyrhythm Architecture}

\subsubsection{The 300Hz LCM}

The Least Common Multiple of 50 and 60 is 300. This means the two standards only synchronize
(return to a common phase point) every 1/300th of a second. At 300\,Hz, both 50\,Hz (6 cycles per
300\,Hz period) and 60\,Hz (5 cycles per 300\,Hz period) divide evenly. A 300\,Hz display, in
principle, could present a frame aligned to either standard's field boundary with zero fractional
offset. This mathematical property underlies the 120\,Hz television standard, which is LCM(24, 25, 30)
and thus accommodates film (24\,fps), PAL video (25\,fps), and NTSC video (30\,fps) without judder.

\subsubsection{3:2 Pulldown}

To transfer 24\,fps film to NTSC's 59.94 fields per second, each film frame is assigned either
2 or 3 video fields in alternating pattern (A=3, B=2, C=3, D=2, \ldots). This ``3:2 pulldown'' causes
a periodic motion cadence artifact visible in slow-motion playback and in certain camera movements.
The pattern repeats every 5 film frames, creating a 4:2 pulldown sequence at the sequence level.
Reverse telecine (``inverse pulldown'') must detect and remove this pattern to recover the original
24\,fps timing for DVD mastering or digital intermediate work.

\subsubsection{Standards Conversion Methods}

\begin{center}
\rowcolors{2}{rowgray}{white}
\begin{tabularx}{\textwidth}{L{3.5cm}XC{2cm}}
\toprule
\rowcolor{secondary}
\tableheader{Method} & \tableheader{Mechanism / Artifacts} & \tableheader{Quality} \\
\midrule
Frame dropping/duplication & Remove or repeat frames; produces temporal stutter & Poor \\
3:2 pulldown & Cadence artifacts; acceptable for film-sourced content & Moderate \\
Field blending & Mix adjacent fields; produces motion blur & Moderate \\
Motion-compensated interpolation & Optical-flow frame synthesis; smooth but requires compute & Good \\
\bottomrule
\end{tabularx}
\end{center}

\subsection{Module E Reference --- High-Frequency Refresh Displays}

\subsubsection{Sample-and-Hold Motion Blur}

Modern flat-panel displays (LCD, OLED) hold each frame on screen for the full frame period, unlike
CRTs which scanned briefly and went dark. The eye, tracking a moving object, integrates the held
image over the frame period, producing motion blur proportional to object velocity $\times$ frame
time. At 60\,Hz (16.7\,ms frame time), this blur is substantial. At 240\,Hz (4.2\,ms), it is
reduced by a factor of four. Formula: blur width (pixels) = motion speed (px/s) $\times$ frame
time (s). At 3840\,px/s motion (1080p screen width per second), 60\,Hz produces 64 pixels of blur;
240\,Hz produces 16 pixels (Medium, Luizbizzio, 2025).

\subsubsection{MPRT vs GtG}

\begin{itemize}[leftmargin=*]
  \item \textbf{GtG (Gray-to-Gray):} Time for a pixel to transition between two luminance levels.
        Relevant for preventing ghosting/trailing.
  \item \textbf{MPRT (Moving Picture Response Time):} Perceived blur duration during eye-tracking
        motion. Combines pixel response time and display persistence. The more perceptually
        relevant metric; MPRT $\leq$ 2\,ms is required for visibly sharp motion on 240+\,Hz displays.
\end{itemize}

\subsubsection{Perceptual Science}

Frontiers in Neuroscience (PMC 8777290, 2022): An EEG study using Steady-State Motion Visual Evoked
Potentials found that increasing refresh rate from 60\,Hz to 120\,Hz significantly improved visual
motion evoked potential amplitude across stimulation frequencies of 7--28\,Hz. For high-velocity
stimuli, 240\,Hz or greater was recommended. This provides peer-reviewed evidence that human motion
perception extracts information beyond the classical ``24\,fps is sufficient'' claim.

\begin{center}
\rowcolors{2}{rowgray}{white}
\begin{tabularx}{\textwidth}{C{2cm}C{2.5cm}C{2.5cm}X}
\toprule
\rowcolor{secondary}
\tableheader{Rate} & \tableheader{Frame Time} & \tableheader{Input Lag vs 60Hz} & \tableheader{Notes} \\
\midrule
60\,Hz & 16.7\,ms & baseline & Standard consumer; NTSC legacy \\
120\,Hz & 8.3\,ms & $-$8.4\,ms & LCM(24/25/30); TV standard \\
144\,Hz & 6.9\,ms & $-$9.8\,ms & Entry gaming; most common upgrade \\
240\,Hz & 4.2\,ms & $-$12.5\,ms & Professional gaming sweet spot \\
360\,Hz & 2.8\,ms & $-$13.9\,ms & Esports; Alienware AW2523HF \\
500\,Hz & 2.0\,ms & est.\ $-$14.7\,ms & Production 2024--2025 \\
1000\,Hz & 1.0\,ms & est.\ $-$15.7\,ms & Research/esports prototype \\
\bottomrule
\end{tabularx}
\end{center}

Input lag reduction from 60\,Hz to 360\,Hz: up to 28\,ms in Fortnite benchmarks at Epic settings
(PCWorld, 2023). Professional esports survey of 1,000+ players: 66\% use 240\,Hz, 11\% use 360\,Hz.

\subsubsection{Variable Refresh Rate and BFI}

VRR standards (VESA Adaptive Sync, AMD FreeSync, NVIDIA G-Sync) synchronize display refresh to GPU
output, eliminating screen tearing and stuttering during frame rate variation. Black Frame Insertion
(BFI) inserts a black frame between content frames, effectively halving persistence without reducing
refresh rate, producing CRT-like clarity. BFI only works effectively at $\geq$240\,Hz to avoid
visible flicker.

\subsection{Module F Reference --- Large Format, 8K, and Scanning Backs}

\subsubsection{Medium Format Sensor Landscape}

\begin{center}
\rowcolors{2}{rowgray}{white}
\begin{tabularx}{\textwidth}{L{3.5cm}C{2cm}C{2.5cm}C{2cm}C{1.5cm}}
\toprule
\rowcolor{secondary}
\tableheader{System} & \tableheader{Resolution} & \tableheader{Sensor Size} & \tableheader{Bit Depth} & \tableheader{Dynamic Range} \\
\midrule
Phase One IQ4 150MP & 150\,MP & 53.7$\times$40.4\,mm & 16-bit & 15\,stops \\
Hasselblad X2D II 100C & 100\,MP & 43.8$\times$32.9\,mm & 16-bit & 15.3\,stops \\
Fujifilm GFX 100S II & 100\,MP & 43.8$\times$32.9\,mm & 14-bit & $\sim$14\,stops \\
Hasselblad H6D-400c MS & 400\,MP (composite) & 53.7$\times$40.4\,mm & 16-bit & 15\,stops \\
Phase One IQ3 100MP & 100\,MP & 53.7$\times$40.4\,mm & 16-bit & 15\,stops \\
\bottomrule
\end{tabularx}
\end{center}

The Phase One sensor (53.7$\times$40.4\,mm) covers approximately the full frame of 645 medium
format film. The Hasselblad/Fuji 43.8$\times$32.9\,mm sensor is a 1.3$\times$ crop of true 645,
fitting between full-frame 35\,mm (36$\times$24\,mm) and full-645. (PetaPixel, Capture Integration
shootout, July 2025; EOSHD medium format guide).

\subsubsection{Trilinear CCD Scanning Backs}

A digital scanning back replaces the matrix area sensor with a single row of trilinear CCD sensors
(three rows: R, G, B) that sweeps mechanically across the focal plane. For each column position, the
sensor captures full red, full green, and full blue data --- no Bayer filter, no color interpolation.
This eliminates the mosaic color estimation that reduces effective resolution in matrix sensors.
Advantages:
\begin{itemize}[leftmargin=*]
  \item True RGB color at every pixel: no Bayer interpolation artifacts
  \item Immune to off-axis light fall-off (no vignetting from lens angle)
  \item Very large scan areas compatible with 4$\times$5 view camera movements (rise, shift, tilt)
  \item File sizes: BetterLight Super8K produces 1.1\,GB 48-bit uncompressed TIFF at maximum resolution
        (12,000$\times$15,990 pixels $\approx$ 192\,MP, 300\,ppi over 40$\times$53 inch scan area)
\end{itemize}
Limitation: exposure time measured in minutes; no moving subjects. Cultural heritage and archival
digitization applications (maps, manuscripts, paintings, architectural drawings).
(Creekside Digital; Wikipedia, Digital Scan Back; Wikipedia, BetterLight)

The 2025 DIY ``Project Gigapixel'' camera (Yannick Richter) achieved 3,200\,MP (80,000$\times$40,000
pixels) using an Epson flatbed scanner's 12-line ILX561K CCD with custom Raspberry Pi control;
a full-resolution capture takes over 30 minutes. (PetaPixel, November 2025)

\subsubsection{8K Standards and NHK Super Hi-Vision}

\begin{center}
\rowcolors{2}{rowgray}{white}
\begin{tabularx}{\textwidth}{L{3cm}C{2.5cm}C{2cm}X}
\toprule
\rowcolor{secondary}
\tableheader{Standard} & \tableheader{Resolution} & \tableheader{Frame Rate} & \tableheader{Notes} \\
\midrule
ITU-R BT.2020 & 7680$\times$4320 & Up to 120\,fps & 10/12-bit; Rec.2020 wide gamut \\
SMPTE ST 2036-1 & 7680$\times$4320 & Up to 120\,fps & UHDTV2; defines MPEG-H / HEVC \\
8K Association & 7680$\times$4320 & 24/30/60\,fps & 600\,nits peak; HEVC; HDMI 2.1 \\
CTA 8K Standard & 7680$\times$4320 & 24/30/60\,fps & LG-aligned; competing with 8KA \\
NHK SHV (BS8K) & 7680$\times$4320 & 60/120\,fps & Launched December 1, 2018 \\
\bottomrule
\end{tabularx}
\end{center}

Total pixel count: 33.2\,million per frame (4$\times$ 4K, 16$\times$ 1080p). NHK's BS8K channel
launched December 1, 2018, carrying documentaries, sports, and remastered films including 2001:
A Space Odyssey (from 70\,mm prints). NHK developed a 133-megapixel 8K CMOS prototype for 60\,fps
capture. The Samsung QN990F (2025) features Mini-LED backlighting, NQ8 AI Gen3 processor for 8K
upscaling, and motion rate up to 240\,Hz. (Wikipedia 8K Resolution; Grokipedia 8K Resolution, 2026)

As of 2026, Samsung holds nearly the entire consumer 8K TV market following LG, Sony, and others
exiting the segment. Dolby Vision is effectively absent from new consumer 8K televisions; Samsung's
ecosystem is HDR10+. (Grokipedia, 2026)

\subsubsection{HDR Standards Comparison}

\begin{center}
\rowcolors{2}{rowgray}{white}
\begin{tabularx}{\textwidth}{L{2.5cm}C{2cm}C{2cm}X}
\toprule
\rowcolor{secondary}
\tableheader{Standard} & \tableheader{Bit Depth} & \tableheader{Metadata} & \tableheader{Notes} \\
\midrule
HDR10 & 10-bit & Static & Open standard; universal baseline \\
HLG & 10-bit & None & BBC/NHK; backward compatible with SDR \\
HDR10+ & 10-bit & Dynamic & Samsung/Amazon; open royalty-free \\
Dolby Vision & 12-bit & Dynamic & Dolby proprietary; exits 8K segment 2026 \\
Dolby Vision 2 & 12-bit & Dynamic & 2025 announcement; no 8K implementations yet \\
\bottomrule
\end{tabularx}
\end{center}

\subsection{Safety and Sensitivity Considerations}

\begin{center}
\rowcolors{2}{rowgray}{white}
\begin{tabularx}{\textwidth}{L{3cm}XC{2cm}}
\toprule
\rowcolor{primary}
\tableheader{Area} & \tableheader{Consideration} & \tableheader{Action} \\
\midrule
Photosensitivity & High-frequency displays and flicker rates may be medically relevant for photosensitive users & ANNOTATE \\
Numerical claims & All frequency values, resolution specs, and timecode mathematics must cite primary standards & GATE \\
Geopolitical framing & The 50/60\,Hz divide has infrastructure-justice implications (Global South legacy adoption) & ANNOTATE \\
Manufacturer claims & Distinguish marketing ``motion rate'' from measured MPRT; flag upscaled vs.\ native 8K & GATE \\
\bottomrule
\end{tabularx}
\end{center}

\subsection{Source Bibliography}

\subsubsection{Standards Bodies and Technical Organizations}

\begin{itemize}[leftmargin=*]
  \item SMPTE ST 12M (SMPTE 12M-1975, revised 1986, split 2008 into ST 12-1 and ST 12-2)
  \item SMPTE ST 12-3 (High Frame Rate Timecode extension)
  \item SMPTE ST 2036-1 (Ultra High Definition Television, UHDTV2)
  \item ITU-R BT.2020 (Parameter Values for Ultra High Definition Television)
  \item ITU-R BT.709 (HDTV standard reference)
  \item IEEE 1588 (Precision Time Protocol)
  \item EBU Technical Recommendation R68 (25\,fps timecode)
  \item VESA DisplayPort 1.4 specification (2016)
  \item HDMI 2.1 specification (HDMI Forum, 2017)
  \item CTA 8K Ultra HD TV standard (2019, LG-aligned)
\end{itemize}

\subsubsection{Peer-Reviewed Research}

\begin{itemize}[leftmargin=*]
  \item Han, Xu, Zheng, et al. (2022). ``Assessing the Effect of the Refresh Rate of a Device on
        Various Motion Stimulation Frequencies Based on Steady-State Motion Visual Evoked Potentials.''
        \textit{Frontiers in Neuroscience}, PMC 8777290.
\end{itemize}

\subsubsection{Manufacturer Documentation}

\begin{itemize}[leftmargin=*]
  \item Hasselblad, ``Medium Format Advantages'' (hasselblad.com/learn/medium-format/)
  \item Hasselblad, ``H6D-400c MS Press Release'' (2018)
  \item Hasselblad, ``X1D II 50C Product Documentation'' (hasselblad.com)
  \item Phase One, IQ3 100MP / IQ4 150MP product specifications
  \item BetterLight, ``Digital Scanning Camera Inserts'' (betterlight.com)
  \item Creekside Digital, ``Large Format Digitization'' (creeksidedigital.com)
  \item Samsung 8K product specifications: QN900D (2024), QN990F (2025)
\end{itemize}

\subsubsection{Reference Articles}

\begin{itemize}[leftmargin=*]
  \item Wikipedia, ``PAL'' (accessed March 2026)
  \item Wikipedia, ``SMPTE timecode'' (accessed March 2026)
  \item Wikipedia, ``8K resolution'' (accessed March 2026)
  \item Wikipedia, ``Digital scan back'' (accessed March 2026)
  \item Grokipedia, ``SMPTE timecode'' (January 2026)
  \item Grokipedia, ``8K resolution'' (March 2026)
  \item Phil Rees, ``SMPTE EBU Timecode'' (philrees.co.uk)
  \item The Broadcast Bridge, ``Timing: Part 4 --- Analog Television'' (March 2022)
  \item PCWorld, ``240Hz is the new 120Hz'' (October 2023)
  \item PetaPixel, ``Photographer Builds 3,200MP Camera Using CCD Scanner Sensor'' (November 2025)
  \item PetaPixel, ``100-Megapixel Camera Image Quality Shootout'' (July 2025)
\end{itemize}

\newpage

%% =====================================================================
%% STAGE 3 — MISSION PACKAGE
%% =====================================================================
\stageheader{STAGE 3 --- MISSION PACKAGE}{tertiary}

\section{Milestone Specification}

\subsection{Mission Objective}

Produce a comprehensive six-module research document on the architecture of visual time standards
covering the 50/60\,Hz power-grid origin, NTSC/PAL engineering, SMPTE timecode mathematics, the
polyrhythm structure at LCM=300\,Hz, high-frequency display science through 1000\,Hz, and large-format
still imaging from medium-format area sensors through scanning backs and 8K acquisition. The final
deliverable is a publication-quality PDF ($\geq$80 pages) with complete citation apparatus, wave-based
parallel production, and a synthesis module tracing the polyrhythm concept across all six domains.

\subsection{Architecture Overview}

\begin{asciibox}
  WAVE 0: Foundation (Sequential, Haiku)
    Source index, shared frequency/resolution schema, module templates
                |
  WAVE 1: Parallel Research Production (Sonnet + Opus, 3 parallel tracks)
    Track A         Track B         Track C
    [A] Temporal    [C] SMPTE       [E] HF Refresh
    [B] NTSC/PAL    [D] Polyrhythm  [F] LF / 8K
                |
  CAPCOM GATE (human review: numerical accuracy, HDR claims, scope check)
                |
  WAVE 2: Synthesis (Sequential, Opus)
    Cross-module polyrhythm through-line
    Integration of LCM mathematics across all six modules
    Safety/sensitivity review
                |
  WAVE 3: Publication + Verification (Sequential, Sonnet)
    Final PDF assembly, bibliography, verification matrix
    CAPCOM final approval
\end{asciibox}

\subsection{System Layers}

\begin{enumerate}[leftmargin=*, label=\textbf{L\arabic*.}]
  \item \textbf{Shared Schema Layer} --- Frequency tables, frame-rate matrices, resolution comparisons
  \item \textbf{Module Research Layer} --- Six parallel documents (one per module)
  \item \textbf{Synthesis Layer} --- Cross-module polyrhythm essay
  \item \textbf{Publication Layer} --- Assembled PDF with complete apparatus
  \item \textbf{Verification Layer} --- Test plan execution, source audit
\end{enumerate}

\subsection{Deliverables}

\begin{center}
\rowcolors{2}{rowgray}{white}
\begin{tabularx}{\textwidth}{C{0.5cm}L{3.5cm}XC{2.5cm}}
\toprule
\rowcolor{tertiary}
\tableheader{\#} & \tableheader{Deliverable} & \tableheader{Acceptance Criteria} & \tableheader{Spec} \\
\midrule
1 & Shared Schema & Frequency/resolution tables, frame-rate conversion matrix & W0 \\
2 & Module A Document & Power grid history, 50/60 origin, 3-standard world & W1-A \\
3 & Module B Document & NTSC/PAL engineering parameters, subcarrier math & W1-A \\
4 & Module C Document & SMPTE timecode, drop-frame proof, all formats & W1-B \\
5 & Module D Document & LCM=300, pulldown, conversion methods & W1-B \\
6 & Module E Document & HF refresh science, MPRT/GtG, VRR, peer-reviewed evidence & W1-C \\
7 & Module F Document & MF sensors, scanning backs, 8K/SHV, HDR fragmentation & W1-C \\
8 & Synthesis Essay & Polyrhythm through-line across all six modules & W2 \\
9 & Final PDF & Assembled document $\geq$80 pages, full bibliography & W3 \\
10 & Verification Report & All 12 success criteria tested, safety checks passed & W3 \\
\bottomrule
\end{tabularx}
\end{center}

\subsection{Component Breakdown}

\begin{center}
\rowcolors{2}{rowgray}{white}
\begin{tabularx}{\textwidth}{L{3cm}L{2cm}C{1.5cm}C{2cm}}
\toprule
\rowcolor{tertiary}
\tableheader{Component} & \tableheader{Dependencies} & \tableheader{Model} & \tableheader{Est.\ Tokens} \\
\midrule
W0: Source index & None & Haiku & 4,000 \\
W0: Shared schema & None & Haiku & 6,000 \\
W1-A: Module A & W0 & Sonnet & 18,000 \\
W1-A: Module B & W0 & Sonnet & 20,000 \\
W1-B: Module C & W0 & Sonnet & 22,000 \\
W1-B: Module D & W0, W1-C partial & Sonnet & 16,000 \\
W1-C: Module E & W0 & Sonnet & 18,000 \\
W1-C: Module F & W0 & Opus & 24,000 \\
W2: Synthesis & All W1 & Opus & 20,000 \\
W3: Assembly & W2 & Sonnet & 12,000 \\
W3: Verification & W3 & Sonnet & 8,000 \\
\midrule
\multicolumn{3}{l}{\textbf{Total estimated}} & \textbf{168,000} \\
\bottomrule
\end{tabularx}
\end{center}

\subsection{Model Assignment Rationale}

\begin{center}
\rowcolors{2}{rowgray}{white}
\begin{tabularx}{\textwidth}{C{1.5cm}C{2cm}XC{2.5cm}}
\toprule
\rowcolor{tertiary}
\tableheader{Model} & \tableheader{Allocation} & \tableheader{Rationale} & \tableheader{Assigned To} \\
\midrule
Opus & $\sim$30\% & Synthesis, cross-module judgment, Module F cultural complexity & W2 synthesis; Module F \\
Sonnet & $\sim$60\% & Module document production, assembly, verification & Modules A-E; W3 \\
Haiku & $\sim$10\% & Shared schemas, scaffolding, source index & W0 \\
\bottomrule
\end{tabularx}
\end{center}

\section{Wave Execution Plan}

\subsection{Wave Summary}

\begin{center}
\rowcolors{2}{rowgray}{white}
\begin{tabularx}{\textwidth}{C{1cm}L{3cm}C{2cm}C{2cm}C{2cm}}
\toprule
\rowcolor{tertiary}
\tableheader{Wave} & \tableheader{Name} & \tableheader{Execution} & \tableheader{Model} & \tableheader{CAPCOM} \\
\midrule
W0 & Foundation & Sequential & Haiku & No \\
W1 & Research Production & Parallel ($\times$3) & Sonnet + Opus & Yes (post) \\
W2 & Synthesis & Sequential & Opus & Yes (pre) \\
W3 & Publication + Verify & Sequential & Sonnet & Yes (final) \\
\bottomrule
\end{tabularx}
\end{center}

\subsection{Wave 0: Foundation}

\begin{center}
\rowcolors{2}{rowgray}{white}
\begin{tabularx}{\textwidth}{L{3cm}XC{2cm}}
\toprule
\rowcolor{tertiary}
\tableheader{Task} & \tableheader{Output} & \tableheader{Model} \\
\midrule
Source index & Structured bibliography JSON; 30+ sources categorized & Haiku \\
Frequency schema & Shared table: 50/60Hz, frame rates, LCM values, line counts & Haiku \\
Resolution schema & Sensor dimensions, pixel counts, bit depths (all modules) & Haiku \\
Module templates & Six document scaffolds with section headers pre-populated & Haiku \\
\bottomrule
\end{tabularx}
\end{center}

\subsection{Wave 1: Parallel Research Tracks}

\textbf{Track A --- Modules A and B} (run simultaneously, no dependencies between them)

\begin{center}
\rowcolors{2}{rowgray}{white}
\begin{tabularx}{\textwidth}{L{2cm}XC{2cm}}
\toprule
\rowcolor{secondary}
\tableheader{Module} & \tableheader{Key Outputs} & \tableheader{Model} \\
\midrule
A & Power grid history; NTSC/PAL/SECAM adoption table; phosphor/flicker physics & Sonnet \\
B & Color subcarrier derivation; 525 vs.\ 625 parameters; PAL phase alternation & Sonnet \\
\bottomrule
\end{tabularx}
\end{center}

\textbf{Track B --- Modules C and D} (Module D may pull Module C output for LCM context)

\begin{center}
\rowcolors{2}{rowgray}{white}
\begin{tabularx}{\textwidth}{L{2cm}XC{2cm}}
\toprule
\rowcolor{secondary}
\tableheader{Module} & \tableheader{Key Outputs} & \tableheader{Model} \\
\midrule
C & 0.1\% offset derivation; drop-frame proof with examples; LTC/VITC/ATC/ST 12-3 & Sonnet \\
D & LCM=300 derivation; 3:2 pulldown walkthrough; conversion method comparison & Sonnet \\
\bottomrule
\end{tabularx}
\end{center}

\textbf{Track C --- Modules E and F} (largest token budget; F uses Opus for sensor complexity)

\begin{center}
\rowcolors{2}{rowgray}{white}
\begin{tabularx}{\textwidth}{L{2cm}XC{2cm}}
\toprule
\rowcolor{secondary}
\tableheader{Module} & \tableheader{Key Outputs} & \tableheader{Model} \\
\midrule
E & MPRT/GtG distinction; sample-and-hold physics; neuroscience citation; display table & Sonnet \\
F & Medium format table; scanning back operation; BetterLight spec; 8K standards table; HDR fragmentation & Opus \\
\bottomrule
\end{tabularx}
\end{center}

\subsection{Wave 2: Synthesis}

\begin{center}
\rowcolors{2}{rowgray}{white}
\begin{tabularx}{\textwidth}{L{4cm}XC{2cm}}
\toprule
\rowcolor{tertiary}
\tableheader{Task} & \tableheader{Description} & \tableheader{Model} \\
\midrule
Polyrhythm through-line & Essay tracing 50/60Hz polyrhythm from power grid through every module & Opus \\
LCM topology & Mathematical through-line: LCM(50,60)=300, LCM(24,25,30)=600, drop-frame as 1/1001 correction & Opus \\
Cross-module consistency & Verify shared terms and numbers match across all six module documents & Opus \\
Sensitivity review & Flag any medical (photosensitivity) or geopolitical framing issues & Opus \\
\bottomrule
\end{tabularx}
\end{center}

\subsection{Wave 3: Publication and Verification}

\begin{center}
\rowcolors{2}{rowgray}{white}
\begin{tabularx}{\textwidth}{L{4cm}XC{2cm}}
\toprule
\rowcolor{tertiary}
\tableheader{Task} & \tableheader{Description} & \tableheader{Model} \\
\midrule
LaTeX assembly & Compile all six module documents + synthesis into master PDF & Sonnet \\
Bibliography final & Consolidate and de-duplicate all citations; verify URLs live & Sonnet \\
Verification matrix & Execute all test plan items; mark pass/fail & Sonnet \\
CAPCOM final & Present to human for review before archive & Human (CAPCOM) \\
\bottomrule
\end{tabularx}
\end{center}

\subsection{Cache Optimization Strategy}

\begin{itemize}[leftmargin=*]
  \item Wave 0 shared schema is placed at the start of every Wave 1 context window (cache prefix)
  \item Frequency/resolution tables computed once in W0 and referenced, not recomputed
  \item Source index JSON is the first attachment in every W1 context to enable consistent citations
  \item 5-minute TTL cache exploited: all three W1 tracks should start within the cache window of W0
  \item Module C's drop-frame proof is a standalone mathematical document; cache as a reusable component
\end{itemize}

\subsection{Token Budget}

\begin{center}
\rowcolors{2}{rowgray}{white}
\begin{tabularx}{\textwidth}{L{3cm}C{2cm}C{2cm}C{2.5cm}}
\toprule
\rowcolor{tertiary}
\tableheader{Wave} & \tableheader{Sessions} & \tableheader{Windows} & \tableheader{Est.\ Tokens} \\
\midrule
W0: Foundation & 1 & 1 & 10,000 \\
W1 Track A & 1 & 2 & 38,000 \\
W1 Track B & 1 & 2 & 38,000 \\
W1 Track C & 1--2 & 2 & 42,000 \\
W2: Synthesis & 1 & 1 & 20,000 \\
W3: Publication & 1 & 1--2 & 20,000 \\
\midrule
\textbf{Total} & \textbf{5--6} & \textbf{9--10} & \textbf{168,000} \\
\bottomrule
\end{tabularx}
\end{center}

\subsection{Dependency Graph}

\begin{asciibox}
  [W0: Source Index + Shared Schema]
        |        |        |
        v        v        v
  [W1-A:  ] [W1-B:  ] [W1-C:  ]     <-- Parallel
  [Mod A  ] [Mod C  ] [Mod E  ]
  [Mod B  ] [Mod D  ] [Mod F  ]
        |        |        |
        v        v        v
  [CAPCOM GATE: numerical + scope review]
                 |
                 v
         [W2: Synthesis]
         [Polyrhythm essay]
         [Cross-module check]
                 |
         [CAPCOM GATE: pre-publication]
                 |
                 v
         [W3: Assembly + Verify]
                 |
         [CAPCOM FINAL]
\end{asciibox}

\section{Test Plan}

\subsection{Test Categories}

\begin{center}
\rowcolors{2}{rowgray}{white}
\begin{tabularx}{\textwidth}{L{3cm}C{2cm}C{2cm}C{2cm}C{2cm}}
\toprule
\rowcolor{tertiary}
\tableheader{Category} & \tableheader{Count} & \tableheader{Priority} & \tableheader{Action} & \tableheader{Target \%} \\
\midrule
Safety-critical & 4 & Mandatory & BLOCK & 15\% \\
Core functionality & 18 & Required & BLOCK & 50\% \\
Integration & 8 & Required & BLOCK & 22\% \\
Edge cases & 5 & Best-effort & LOG & 13\% \\
\midrule
\textbf{Total} & \textbf{35} & & & 100\% \\
\bottomrule
\end{tabularx}
\end{center}

\subsection{Safety-Critical Tests}

\begin{center}
\rowcolors{2}{rowgray}{white}
\begin{tabularx}{\textwidth}{L{1.5cm}L{3cm}X}
\toprule
\rowcolor{tertiary}
\tableheader{Test ID} & \tableheader{Verifies} & \tableheader{Expected Behavior} \\
\midrule
SC-SRC & Source quality & All citations traceable to standards bodies, peer-reviewed journals, or manufacturer documentation. Zero entertainment media. \\
SC-NUM & Numerical attribution & Every frequency value, frame rate, resolution, and timecode formula attributed to a specific standard or source \\
SC-ADV & No policy advocacy & Document presents technical findings without advocating for one standard over another \\
SC-MED & Photosensitivity disclaimer & Any discussion of flicker rates includes appropriate note re: photosensitive readers \\
\bottomrule
\end{tabularx}
\end{center}

\subsection{Core Functionality Tests (selected)}

\begin{center}
\rowcolors{2}{rowgray}{white}
\begin{tabularx}{\textwidth}{L{1.5cm}L{3.5cm}X}
\toprule
\rowcolor{tertiary}
\tableheader{Test ID} & \tableheader{Verifies} & \tableheader{Expected Behavior} \\
\midrule
CF-01 & Power grid causal chain (SC-1) & Module A contains explicit source linking 50/60Hz power frequency to broadcast field rate \\
CF-02 & Three-standard table (SC-1) & NTSC/PAL/SECAM table with regions, line counts, field rates, color systems present \\
CF-03 & NTSC subcarrier derivation (SC-2) & Module B contains exact value 315/88 MHz = 3.57954545 MHz with derivation steps \\
CF-04 & PAL phase alternation (SC-2) & Module B explains the 1-line delay correction and Hanover bars suppression \\
CF-05 & 0.1\% offset origin (SC-3) & Module C traces offset to color subcarrier compatibility requirement with exact factor 1000/1001 \\
CF-06 & Drop-frame arithmetic (SC-3) & Module C shows 18 frames per 10 min block = 18,000 at 30fps, residual 86.4ms/day \\
CF-07 & All timecode formats (SC-4) & Module C documents LTC, VITC, ATC, ST 12-3 with standards references \\
CF-08 & LCM(50,60)=300 proof (SC-5) & Module D contains explicit LCM derivation and explanation of sync point frequency \\
CF-09 & 3:2 pulldown walkthrough (SC-6) & Module D shows the A=3/B=2/C=3/D=2 field assignment pattern with example \\
CF-10 & Neuroscience citation (SC-7) & Module E cites PMC 8777290 with correct findings ($\geq$120Hz recommendation) \\
CF-11 & MPRT vs GtG distinction (SC-8) & Module E provides separate definitions with example values \\
CF-12 & VRR and BFI operation (SC-8) & Module E describes both technologies with operational conditions \\
CF-13 & Sensor dimension table (SC-9) & Module F table covers Phase One/Hasselblad/Fuji with sensor dimensions in mm \\
CF-14 & Scanning back operation (SC-10) & Module F explains trilinear CCD sweep, no interpolation, scan time \\
CF-15 & BetterLight spec data (SC-10) & Module F contains 12,000x15,990, 1.1GB, 48-bit, 300ppi values with source \\
CF-16 & 8K standards table (SC-11) & Module F table covers ITU-R BT.2020, ST 2036-1, 8KA, NHK SHV \\
CF-17 & HDR fragmentation (SC-11) & Module F documents Samsung near-monopoly and Dolby Vision exit from 8K \\
CF-18 & Synthesis polyrhythm (SC-12) & Synthesis module traces polyrhythm concept through all six domain areas \\
\bottomrule
\end{tabularx}
\end{center}

\subsection{Integration Tests (selected)}

\begin{center}
\rowcolors{2}{rowgray}{white}
\begin{tabularx}{\textwidth}{L{1.5cm}L{4cm}X}
\toprule
\rowcolor{tertiary}
\tableheader{Test ID} & \tableheader{Verifies} & \tableheader{Expected Behavior} \\
\midrule
IN-01 & Mod A $\rightarrow$ B cascade & Module B's NTSC subcarrier derivation explicitly references the 60Hz power-grid standard established in Module A \\
IN-02 & Mod B $\rightarrow$ C cascade & Module C's 0.1\% offset derives from the same color subcarrier value established in Module B \\
IN-03 & Mod C $\rightarrow$ D cascade & Module D's LCM analysis references the 29.97fps drop-frame context from Module C \\
IN-04 & Mod D $\rightarrow$ E cascade & Module E's 120Hz as LCM(24,25,30) references the polyrhythm mathematical framework of Module D \\
IN-05 & Mod E $\rightarrow$ F bridge & Module F's 8K frame rates (24/30/60fps) are contextualized against Module E's refresh science \\
IN-06 & Cross-module consistency & The value 29.97fps appears consistently as 30/1.001 across Modules B, C, and D \\
IN-07 & Self-containment & A reader without prior knowledge can follow the full power-grid-to-scanning-back lineage from Module A through F \\
IN-08 & Synthesis integration & Synthesis essay references specific technical findings from all six modules, not generic claims \\
\bottomrule
\end{tabularx}
\end{center}

\subsection{Verification Matrix}

\begin{center}
\rowcolors{2}{rowgray}{white}
\begin{tabularx}{\textwidth}{XL{3.5cm}C{1.5cm}}
\toprule
\rowcolor{tertiary}
\tableheader{Success Criterion} & \tableheader{Test ID(s)} & \tableheader{Status} \\
\midrule
SC-1: Power grid causal chain & CF-01, CF-02, IN-01 & Pending \\
SC-2: NTSC/PAL engineering parameters & CF-03, CF-04 & Pending \\
SC-3: 0.1\% offset and drop-frame & CF-05, CF-06, IN-02 & Pending \\
SC-4: All SMPTE timecode formats & CF-07 & Pending \\
SC-5: LCM=300 polyrhythm foundation & CF-08, IN-03 & Pending \\
SC-6: Standards conversion methods & CF-09 & Pending \\
SC-7: Peer-reviewed neuroscience citation & CF-10 & Pending \\
SC-8: MPRT/GtG/VRR/BFI documentation & CF-11, CF-12, IN-04 & Pending \\
SC-9: Medium format sensor landscape & CF-13, IN-05 & Pending \\
SC-10: Scanning back operation and specs & CF-14, CF-15 & Pending \\
SC-11: 8K standards and HDR fragmentation & CF-16, CF-17 & Pending \\
SC-12: Polyrhythm synthesis across all modules & CF-18, IN-08, IN-07 & Pending \\
\bottomrule
\end{tabularx}
\end{center}

\section{Mission Crew Manifest}

\begin{center}
\rowcolors{2}{rowgray}{white}
\begin{tabularx}{\textwidth}{L{2cm}L{3cm}X}
\toprule
\rowcolor{tertiary}
\tableheader{Role} & \tableheader{Agent} & \tableheader{Responsibility} \\
\midrule
FLIGHT & Orchestrator & Mission direction; go/no-go at CAPCOM gates; wave sequencing \\
PLAN & Planner & Wave decomposition; parallel track scheduling; dependency graph management \\
TOPO & Topology & Agent team structure; mid-mission restructuring if a track stalls \\
SCOUT & Research & Gap-fill web research; source validation; secondary searches \\
ARCH & Architecture & Cross-cutting structural decisions; schema consistency \\
EXEC-A & Executor Alpha & Wave 1 Track A production (Modules A and B) \\
EXEC-B & Executor Beta & Wave 1 Track B production (Modules C and D) \\
EXEC-C & Executor Gamma & Wave 1 Track C production (Modules E and F) \\
CRAFT-SIGNAL & Domain Specialist & Broadcast engineering, SMPTE, sync chains, display science \\
CRAFT-OPTICS & Domain Specialist & Large-format photography, sensor architecture, 8K imaging \\
INTEG & Integration & Cross-module relationship validation; polyrhythm thread \\
VERIFY & Verification & Source validation; numerical accuracy; test plan execution \\
SECURE & Safety & SC-SRC, SC-NUM, SC-MED enforcement; photosensitivity flags \\
ANALYST & Pattern Analyst & LCM topology detection; cross-module mathematical consistency \\
CAPCOM & HITL Interface & Human approval channel; three CAPCOM gates (post W1, pre W2, post W3) \\
BUDGET & Resource & Token budget tracking per wave; session planning \\
RETRO & Retrospective & Post-wave analysis; schema version control \\
SURGEON & Health Monitor & Research quality degradation detection; source drift alerts \\
LOG & Chronicle & Commit messages; milestone summaries; audit trail \\
\bottomrule
\end{tabularx}
\end{center}

\section{Execution Summary}

\begin{center}
\rowcolors{2}{rowgray}{white}
\begin{tabularx}{\textwidth}{L{4cm}X}
\toprule
\rowcolor{tertiary}
\tableheader{Metric} & \tableheader{Value} \\
\midrule
Mission name & The Polyrhythm Standard \\
Activation profile & Fleet (6+ modules) \\
Module count & 6 (A: Origins, B: NTSC/PAL, C: SMPTE, D: Polyrhythm, E: HF Display, F: LF / 8K) \\
Crew roles & 19 \\
Wave count & 4 (W0--W3) \\
Parallel tracks & 3 (Wave 1) \\
CAPCOM gates & 3 \\
Total estimated tokens & 168,000 \\
Estimated sessions & 5--6 \\
Estimated context windows & 9--10 \\
Primary model split & Opus 30\% / Sonnet 60\% / Haiku 10\% \\
Safety-critical tests & 4 (all BLOCK) \\
Total tests & 35 \\
Output document & Publication-quality PDF $\geq$80 pages \\
\bottomrule
\end{tabularx}
\end{center}

%% CLOSING QUOTE
\vspace{2em}
\begin{quotebox}
\textit{``The number 50 and the number 60 share nothing below 300. The history of video engineering
is the story of what happens in those spaces between --- the standards conversion equipment, the
drop-frame hack, the 120\,Hz LCM display, the scanning back's trilinear patience. Meaning lives
not in the frequencies themselves but in the architecture built to reconcile them.''}

\vspace{0.5em}
\hfill{\small\sffamily\textcolor{subtlegray}{--- Research Mission Vision, GSD Ecosystem, March 2026}}
\end{quotebox}

\end{document}
